\section{Targeted Vulnerabilities and Risks}
\label{sec:landscape}

\paragraph{Memory-safety defects in C network services.}
nginx is written in C, and it parses attacker-controlled input at the network
edge. Its history contains exactly the defects that this implies. Examples
are a stack buffer overflow that is reachable through chunked transfer
encoding~\cite{cve2013:nginx}, an off-by-one heap write in the DNS resolver
that was remotely triggerable across a twelve-year span of
versions~\cite{cve2021:resolver}, and repeated memory-disclosure defects in
optional modules that many operators did not know they had compiled
in~\cite{cve2022:mp4}. These are the defects that the component-level analysis of Section~\ref{sec:thrust3} targets, and they are why that analysis cannot be run against production.

\paragraph{Configuration-induced vulnerabilities.}
The larger practical risk is not a flaw in nginx, but a flaw in the way that
nginx was configured. The directive language admits well-documented footguns
whose consequences are severe and whose presence is invisible to the
operator. For example, an \texttt{alias} directive without a trailing slash
yields path traversal, disabling slash merging in front of an upstream that
normalizes differently yields an access-control bypass, a
\texttt{proxy\_pass} directive that uses a variable silently disables URI
normalization, and parsing mismatches between nginx and its upstream enable
request smuggling~\cite{cve2019:smuggling}. This class of defect is systemic
rather than incidental, and it has been shown to generalize across proxies
and application servers~\cite{orangetsai}. Crucially, these vulnerabilities exist only in the \emph{deployed configuration}, and cannot be found by
analyzing the nginx source alone, nor in a staging environment that does not reproduce production. They are the paradigmatic case for a fidelity-controlled twin.

\paragraph{Cross-component chains.}
The findings that matter most are compositions: a permissive proxy rule that
exposes an internal endpoint, combined with a deserialization flaw in a
client library, combined with a credential that is reachable from the
compromised service. Each link in such a chain may be individually
unremarkable, but the chain as a whole is critical.

\paragraph{Socio-technical risks.}
Deployments do not install nginx. They install a container image, a distribution package, or an ingress
controller, such as \texttt{ingress-nginx}, that contains some build of nginx, with some set of modules, at some version, and from some registry. The
Log4Shell incident showed how a defect in a transitive dependency propagates
instantly across every system that unknowingly embeds it~\cite{log4shell}.
\sysname records provenance and Software Bills of Materials in the knowledge
base, as described in Section~\ref{sec:thrust1}, so that the analysis can
reason about attack paths that originate in a dependency which the operator
never chose explicitly. In doing so, we follow the recommended practices of
CISA and NSA~\cite{cisansa} and the criteria of the OpenSSF~\cite{openssf}.
Two further risks are in scope. First, the \emph{third-party module}
ecosystem of nginx is maintained outside the core project and receives a
fraction of its scrutiny, yet modules are compiled into the same address
space with the same privileges, and an unmaintained module is thus an
unpatched vulnerability that has no advisory to trigger an update. Second,
the concentration of \emph{maintainer and governance} authority is itself a
security property, as the xz-utils backdoor showed that a patient adversary
can obtain maintainer trust and introduce a backdoor that survives ordinary
review~\cite{xzbackdoor}.
%And \emph{insider and
%poisoned-component} risk is covered because the Red Team Agents plan from
%arbitrary initial footholds, so an adversary holding valid credentials or
%controlling a build artifact is expressible as a starting condition in the
%same planning formulation.

\paragraph{Prior incidents and threat model.}
These classes of defect are not hypothetical, and the incidents that
instantiate them share a structure that motivates the design of this project.
The clearest example is the set of defects that were disclosed in 2025 in the
\texttt{ingress-nginx} admission controller, collectively named
``IngressNightmare,'' which permitted unauthenticated remote code execution
and the disclosure of cluster-wide secrets on a large fraction of Kubernetes
deployments~\cite{ingressnightmare}. The relevant point is not the defect itself, but its reachability, as whether a given cluster was exploitable depended on the network policy, on whether the controller was reachable from the workload pods, and on the surrounding configuration. That question is answerable only against a faithful replica, and it is exactly the question no operator could safely ask of production. The same structure recurs in the resolver defect of 2021, whose
exploitability depended on whether an attacker could influence DNS responses
from the network position of the deployment~\cite{cve2021:resolver}, and in
request smuggling, which depends on the specific pairing of proxy and
upstream~\cite{cve2019:smuggling,orangetsai}. In every case, the advisory
establishes that a defect exists, and only the deployment determines whether
it matters.

Accordingly, we consider an adversary whose goal is unauthorized access to a
designated high-value asset, such as a patient database, a credential store,
or a cluster control plane. We model four initial positions: an
\emph{external} adversary, who has network access only to internet-facing
services; a \emph{foothold} adversary, who has compromised one low-privilege
component; an \emph{insider} adversary, who holds valid credentials for one
role; and a \emph{supply-chain} adversary, who controls one dependency or
build artifact. The adversary is assumed to know the open-source composition
of the infrastructure, since that information is public, but not its
deployment-specific configuration or its secrets. Two assumptions bound the
work. First, \sysname evaluates the security of the \emph{deployment}, and physical
attacks, social engineering, and the compromise of the virtualization
platform are out of scope. Second, \sysname never attacks the production
system, and any conclusion about production is therefore bounded by the
measured fidelity of the twin, which Section~\ref{sec:thrust2} makes
explicit and quantitative rather than assumed.
